\chapter{SİBER TEHDİT İSTİHBARATI VE THREAT HUNTING}

\section{Cyber Threat Intelligence (CTI) Fundamentals}

\subsection{Threat Intelligence Lifecycle ve Collection Methods}
Tehdit istihbaratı yaşam döngüsü, ham veriyi değerli istihbarata dönüştüren ve sürekli bir geri bildirim döngüsü ile kendini yenileyen, yapılandırılmış bir süreçtir. Bu döngü, CTI programının temelini oluşturur ve reaktif bir modelden proaktif bir savunma yaklaşımına geçişin temelini atar.

\subsubsection{Yaşam Döngüsü Aşamaları}
\begin{enumerate}
    \item \textbf{Gereksinimler (Requirements):} (PIR) Kritik varlıkları ve riskleri tanımlama.
    \item \textbf{Toplama (Collection):} OSINT, Ticari Beslemeler, ISAC'ler ve Karanlık Web izleme dahil olmak üzere geniş yelpazede ham veri toplama.
    \item \textbf{İşleme (Processing):} Veri normalizasyonu, tekilleştirme ve zenginleştirme.
    \item \textbf{Analiz ve Yorumlama (Analysis and Interpretation):} Veriden eyleme dönüştürülebilir içgörüler çıkarma.
    \item \textbf{Yayma (Dissemination):} İstihbaratı ilgili paydaşlara (CISO, SOC Analist) uygun formatta sunma.
    \item \textbf{Geri Bildirim (Feedback):} Sürecin etkinliğini artırmak için sürekli geri bildirim toplama.
\end{enumerate}

\subsubsection{CTI Yaşam Döngüsü Aşamaları ve Uygulamaları}
\begin{tabular}{|p{4cm}|p{6cm}|p{4cm}|}
\hline
\hline
\textbf{Aşama} & \textbf{Kısa Tanım} & \textbf{Hedef}  \\
\hline
\textbf{Gereksinimler} & İhtiyaç duyulan istihbaratın belirlenmesi & Kaynakları kritik risklere odaklamak  \\
\hline
\textbf{Toplama} & Ham verinin kaynaklardan alınması & Gereksinim odaklı veri toplama  \\
\hline
\textbf{İşleme} & Ham veriyi analiz formatına getirme & Veri yapılandırma ve zenginleştirme  \\
\hline
\textbf{Analiz} & Veriden içgörüler çıkarma & Tehdit bağlamsallaştırma  \\
\hline
\textbf{Yayma} & İstihbaratı paydaşlara iletme & Doğru kitleye doğru formatta sunum  \\
\hline
\textbf{Geri Bildirim} & Etkinlik geri bildirimi toplama & Sürekli iyileştirme  \\
\hline
\hline
\end{tabular}


\subsection{Strategic, Tactical, Technical ve Operational Intelligence}
Siber tehdit istihbaratı, hedef kitlenin ihtiyaçlarına göre dört ana kategoriye ayrılır.

\subsubsection{Stratejik İstihbarat}
Genel tehdit ortamına ilişkin üst düzey bir bakış açısı sunar. Güvenlik yatırımları ve kurumsal politikalar gibi uzun vadeli stratejik kararlara rehberlik eder. (\textbf{Hedef Kitle:} C-seviyesi yöneticiler)

\subsubsection{Operasyonel İstihbarat}
Belirli bir tehdit aktörünün TTP'lerini, motivasyonlarını ve altyapısını anlamayı sağlar. Proaktif tehdit avcılığı ve olay müdahale planlaması için temel oluşturur. (\textbf{Hedef Kitle:} IR ekipleri, tehdit avcıları)

\subsubsection{Taktik İstihbarat}
Ağ ve uç noktalarda tehditleri tespit etmeye yardımcı olan kısa ömürlü, teknik göstergelerden (IOC'ler) oluşur. Otomatik tehdit tespiti için güvenlik kontrollerine entegre edilir.

\subsubsection{Teknik İstihbarat}
Taktik istihbaratın daha derin teknik detaylarını içerir; kötü amaçlı yazılım analizi ve tersine mühendislik ile yeni imza tabanlı tespit kuralları oluşturmayı amaçlar.

\subsubsection{İstihbarat Türleri Karşılaştırma Tablosu}
\begin{tabular}{|p{4cm}|p{6cm}|p{4cm}|}
\hline
\hline
\textbf{Özellik} & \textbf{Stratejik} & \textbf{Operasyonel}  \\
\hline
\textbf{Hedef Kitle} & C-seviyesi yöneticiler & IR ekipleri, tehdit avcıları  \\
\hline
\textbf{Odak Noktası} & Genel eğilimler, risk yönetimi & Saldırı kampanyaları, TTP'ler  \\
\hline
\textbf{Zaman} & Uzun vadeli (aylar/yıllar) & Orta vadeli (hafta/ay)  \\
\hline
\textbf{Amaç} & Politika ve yatırım belirleme & Müdahale ve proaktif avcılık  \\
\hline
\hline
\end{tabular}

\subsection{Threat Actor Profiling ve Attribution Challenges}

\subsubsection{Tehdit Aktörü Profilleme}
Saldırının arkasındaki tehdit aktörünün kimliğini, hedeflerini, TTP'lerini, motivasyonlarını ve kullandığı altyapıyı kapsamlı bir şekilde analiz eder.

\subsubsection{Atıf Zorlukları (Attribution Challenges)}
Bir siber saldırıyı kesin olarak belirli bir tehdit aktörüne atfetmek (attribution), Gizleme (Obfüskasyon), Yanlış Bayrak Operasyonları ve Araç/TTP Paylaşımı nedeniyle karmaşıktır. Bu nedenle birincil hedef, kimin saldırdığından ziyade, \textbf{davranışsal kanıtlara (TTP'ler)} dayalı savunma stratejileri geliştirmek olmalıdır.

\subsection{Diamond Model ve Kill Chain Analysis}
Siber saldırıları anlamak ve analiz etmek için kullanılan iki önemli analitik çerçeve, Lockheed Martin'in Siber Kill Chain’i ve MITRE’nin Saldırı Analizinin Elmas Modeli’dir (Diamond Model of Intrusion Analysis).

\subsubsection{Siber Kill Chain (Siber Saldırı Zinciri)}
Yedi aşamalı, doğrusal bir saldırı sürecini sunar: Keşif, Silahlandırma, Teslimat, İstismar, Kurulum, Komuta ve Kontrol, ve Hedefler Üzerindeki Eylemler. Olay müdahalesi için adım adım bir yol haritası sunar.

\subsubsection{Diamond Model of Intrusion Analysis}
Saldırıları dört temel bileşen arasındaki ilişkilere odaklanarak inceler: \textbf{Saldırgan (Adversary), Kabiliyet (Capability), Altyapı (Infrastructure)} ve \textbf{Kurban (Victim)}. Doğrusal değildir; tehdit istihbaratı ve proaktif tehdit avcılığı için daha kullanışlıdır, saldırının \textit{neden} ve \textit{kim tarafından} yapıldığına dair bütüncül bir resim sunar.

\subsubsection{Kill Chain ve Diamond Model Karşılaştırma Tablosu}
\begin{tabular}{|p{4cm}|p{6cm}|p{4cm}|}
\hline
\hline
\textbf{Özellik} & \textbf{Siber Kill Chain} & \textbf{Diamond Model} \\
\hline
\textbf{Odak} & Saldırı aşamaları ve süreci & Aktör-kabiliyet-altyapı-kurban ilişkileri \\
\hline
\textbf{Granülarite} & Spesifik ve doğrusal & Geniş, bağlamsal ve ilişkisel \\
\hline
\textbf{Uygulama} & Olay müdahalesi, taktiksel savunma & Tehdit istihbaratı, proaktif avcılık \\
\hline
\textbf{Avantajlar} & Adım adım müdahale yol haritası & Motivasyon ve kabiliyet analizi \\
\hline
\hline
\end{tabular}

\subsection{Intelligence Requirements ve Priority Intelligence Requirements (PIR)}
Priority Intelligence Requirements (PIR'lar), kuruluşun misyonu ve hedefleri açısından tehdit ortamında en kritik konuları temsil eder. CTI programı için kaynak tahsisini ve analistlerin stratejik hedeflere odaklanmasını sağlar.

\subsubsection{PIR'ların Kritik Öğeleri}
Açık ve Spesifik Odak, Ölçülebilir Hedefler, Kuruluşa Alaka Düzeyi, Esneklik ve Adaptasyon.

\subsubsection{PIR'dan Tehdit Avcılığına Dönüşümün Teknik Mekaniği}
PIR'lar $\rightarrow$ Tematik Anahtar Kelimeler $\rightarrow$ Hedeflenen Tehdit Aktörlerinin TTP'leri $\rightarrow$ Tehdit Avcılığı Senaryoları. Bu, avcılık senaryolarının rastgele keşifler yerine, en yüksek riske sahip TTP'ler üzerine odaklanmasını sağlar.

\section{Threat Intelligence Platforms ve Standards}

\subsection{MITRE ATT\&CK Framework Integration}
ATT\&CK, saldırgan davranışlarını sistematik olarak kataloglayan küresel bir bilgi tabanıdır. Kuruluşların savunma yeteneklerini Taktikler, Teknikler ve Alt Teknikler ile haritalandırmasına olanak tanır.

\subsubsection{SIEM Tespit Kuralı Eşleme Metodolojisi}
Tespit kurallarını (\textit{Detection Rules}) ATT\&CK taktikleri ve teknikleri ile ilişkilendirme. Kural tanımına ATT\&CK ID'lerinin öznitelik (\textit{attribute}) olarak eklenmesi, savunma kapsamı metriklerinin hesaplanmasını sağlar.

\subsection{STIX/TAXII Standards ve Information Sharing}
STIX (Structured Threat Information eXpression) ve TAXII (Trusted Automated eXchange of Intelligence Information), siber tehdit istihbaratının yapısal ve güvenli değişimini sağlayan standartlardır.

\subsubsection{STIX 2.1 Veri Modelinin Teknik Detayları}
Tehdit istihbaratını (Malware, Indicators) tanımlayan \textbf{STIX Domain Objects (SDO)}, gözlemlenebilir verileri (IP, hash) tanımlayan \textbf{STIX Cyber Observable Objects (SCO)} ve bunlar arasındaki semantik bağlantıyı kuran \textbf{STIX Relationship Objects (SRO)} kullanılır.

\subsubsection{TAXII 2.1 İstemci Entegrasyonu ve Veri Çekme}
TAXII, STIX verilerinin güvenli, otomatik ve talebe dayalı (RESTful API) değişimini sağlayan protokoldür. İstemciler, büyük veri hacimlerini bütünlük içinde işlemek için sayfalama (\textit{pagination}) mekanizmalarını doğru yönetmelidir.

\subsection{Threat Intelligence Platform (TIP) Selection ve Implementation}
TIP, ham veriyi işleyen, zenginleştiren, önceliklendiren ve savunma sistemlerine yayan merkezi bir sistemdir.

\subsubsection{TIP Mimari Entegrasyon Kriterleri}
Platformun Toplama, İşleme/Dönüşüm, Zenginleştirme ve Dağıtım fonksiyonlarını desteklemesi ve mevcut güvenlik mimarisiyle entegre olması zorunludur (örneğin OAUTH2 Token, Verify SSL kullanımı).

\subsubsection{CTI Veri Dönüşümü ve Standardizasyon Boru Hatları}
Ham verinin, analitik süreçlerin gereksinimlerine uygun, tutarlı ve anlamlı bir formata dönüştürülmesi. Kritik adım, dahili verinin harici istihbarat (itibar skoru, SSL analizi) ile \textbf{Zenginleştirilmesidir} (\textit{Enrichment}).

\subsection{Indicators of Compromise (IOC) Management}
IOC'ler, bir ihlalin kanıtı olan atomik veri parçalarıdır. IOC yönetimi, bunların hızlıca toplanmasını, doğrulanmasını ve savunma mekanizmalarına enjekte edilmesini içerir.

\subsubsection{IOC Yaşam Döngüsü ve Önceliklendirme Metodolojisi}
IOC'lerin güvenilirliği (\textbf{STIX Confidence}) ve potansiyel Yanlış Pozitif (\textbf{FP}) riski ile önceliklendirilmesi. Yaşam döngüsü, IoC'lerin \textbf{eskime} (\textit{expiry/deprecation}) fazlarını da içermelidir.

\subsubsection{Yüksek Hacimli IOC Engellemesi}
Binlerce IP adresini hızlıca bloklamak için Linux'ta \\texttt{ipset} ve \\texttt{iptables} entegrasyonu kullanılır.

\subsection{Tactics, Techniques, and Procedures (TTP) Analysis}
TTP analizi, saldırganın davranış kalıplarını ve yöntemlerini kapsar; davranışsal ve proaktif tehdit avcılığına geçişi sağlar.

\subsubsection{ATT\&CK TTP'lerine Dayalı Hipotez Geliştirme}
Tehdit avcılığı hipotezleri, spesifik bir saldırgan davranışına odaklanmalıdır (örneğin, "X tehdit grubu T1059.001 PowerShell tekniğini kullanarak kalıcılık sağlar").

\subsubsection{PowerShell Evasion TTP'lerinin Tespiti ve Analizi}
Saldırganlar tarafından kullanılan \textbf{-encodedcommand} gibi gizleme tekniklerinin tespiti için \textbf{Windows Event ID 4104 (Scriptblock Logging)} loglarının incelenmesi kritik öneme sahiptir.

\section{Threat Hunting Metodolojileri ve Teknikleri}
Tehdit avcılığı, geleneksel güvenlik sistemlerinin atladığı tehditleri proaktif olarak arama pratiğidir.

\subsection{Hipotez Odaklı Tehdit Avcılığı Yaklaşımları}
Analistin CTI, TTP'ler veya anomali uyarılarına dayanarak formüle ettiği bilinçli bir varsayımla başlar. KQL, SPL gibi sorgu dilleriyle hipotezlerin test edilmesi temel adımdır.

\subsection{Data-driven Hunting ve Statistical Analysis}
Önceden belirlenmiş bir hipotez yerine, geniş veri kümelerinde \textbf{yığın sayımı} (\textit{stack counting}) ve zaman serisi analizi ile anormallikler aranır. Nadir ve yüksek riskli olayları yüzeye çıkarmaya odaklanır.

\subsection{Behavioral Analytics ve Anomaly Detection}
Kullanıcıların ve varlıkların "normal" aktivite modellerini öğrenerek, taban çizgiden küçük sapmaları tespit etmeye odaklanır (\textbf{UEBA - User and Entity Behavior Analytics}). Anormal süreç yaratımı ve anormal ağ bağlantıları izlenir.

\subsection{Hunt Team Organization ve Skill Development}
Avcılık ekipleri, Host Analizi, Ağ Trafik Analizi, Adli Bilişim, Bellek Analizi ve Tersine Mühendislik gibi alanlarda derinlemesine teknik uzmanlığa sahip olmalıdır. \textbf{Volatility Framework} bellek analizi için standart araçtır.

\subsection{Threat Hunting Metrics ve Success Measurement}
Avcılık faaliyetinin etkinliğini ölçmek için \textbf{Dwell Time} (Gizlenme Süresi), \textbf{Breakout Time} ve \textbf{ATT\&CK TTP Kapsam Metrikleri} kullanılır. Başarı, manuel avlarla proaktif olarak tespit edilen olayların, reaktif olaylara oranlanmasıyla ölçülür.

\section{Advanced Persistent Threat (APT) Detection}
Gelişmiş saldırganların gizlenme ve kalıcılık tekniklerine odaklanır.

\subsection{APT Lifecycle ve Long-term Persistence Techniques}

\subsubsection{Zamanlanmış Görev ve WMI Kalıcılığı Tespiti}
APT'ler kalıcılık için \textbf{Zamanlanmış Görevleri} (T1053.005) ve \textbf{WMI Olay Aboneliklerini} (T1546.003) kullanır. Sysmon Event ID 19, 20 ve 21 ile WMI kalıcılığı zincirleme olay olarak izlenir.

\subsection{Lateral Movement Detection ve Analysis}

\subsubsection{PsExec ve Uzaktan Hizmet Yürütme Tespiti}
PsExec kullanımı, \texttt{\textbackslash PSEXESVC.exe} dosya oluşumu (Sysmon EID 11) ve \textbf{Adlandırılmış Boru (\textit{Named Pipe})} bağlantı olaylarının (Sysmon EID 17/18) korelasyonu ile tespit edilir.

\subsubsection{Kerberos Golden Ticket Saldırısı Tespiti}
Domain Controller loglarında \textbf{Event ID 4769} ile anormal hata kodları (\texttt{0x6}) ve aşırı uzun ömürlü (10 yıldan fazla) bilet ömürlerinin tespiti.

\subsection{Living-off-the-Land Techniques Identification}

\subsubsection{Yerel İkili Dosya (LOLBins) Kötüye Kullanımında Görünürlük}
\texttt{PowerShell}, \texttt{CertUtil} ve \texttt{Mshta.exe} gibi meşru yerleşik araçların kötüye kullanımı, yalnızca tam komut satırı argümanlarının yakalanmasıyla (Sysmon EID 1) tespit edilebilir.

\subsubsection{Gelişmiş PowerShell Betik Bloğu Günlüğü Analizi}
APT'nin kod gizleme (\textit{obfuscation}) tekniğini kırmak için \textbf{Windows Event ID 4104 (Script Block Logging)} ile \texttt{IEX}, \texttt{FromBase64String} gibi deobfüskasyon anahtar kelimeleri aranır.

\subsection{Command and Control (C2) Communication Analysis}

\subsubsection{C2 İşaretleşme (Beaconing) Davranışsal Tespiti}
Trafik içeriğinden ziyade, bağlantı başlangıç zamanları arasındaki tutarlı aralıkların (\textbf{Düşük Jitter}) analizi ve tutarlı paket boyutlarının istatistiksel tespiti.

\subsubsection{DNS Tünelleme Tespiti}
DNS sorgularının entropi ve QNAME uzunluklarının hesaplanması, nadir DNS kayıt tiplerinin (TXT, NULL) anormal kullanımı.

\subsection{APT Attribution ve Threat Group Tracking}

\subsubsection{Atıf Metodolojisinin Temelleri ve Amacı}
Saldırganın motivasyonunu, yeteneğini ve operasyonel amacını anlamayı sağlar. Savunma kaynaklarının önceliklendirilmesini destekler.

\subsubsection{TTP Analizi ve Korelasyonu ile Atıf Adımları}
Teknik eserlerin ATT\&CK'e haritalanması, TTP'lerin kronolojik olarak sıralanması (\textbf{TTP Sequencing}) ve bilinen grupların altyapı/kod tekrar kullanımı ile çapraz referanslama.

\section{Malware Analysis ve Reverse Engineering}

\subsection{Static Malware Analysis Techniques ve Tools}
Dosya başlıkları, PE/ELF yapıları ve dize (strings) analizi ile zararlı yazılımın amacı, işlevselliği ve hedeflediği API'ler hakkında ilk bilgilerin toplanması. Yüksek entropi, paketlenmiş (packed) kodu işaret eder.

\subsection{Dynamic Analysis ve Sandbox Environments}
Kötü amaçlı yazılımın kontrollü bir ortamda yürütülmesi ve çalışma zamanı davranışlarının izlenmesi. \textbf{Volatility Framework}, bellek dökümü analizi için kullanılır. \textbf{netscan} ile C2 tespiti, \textbf{malfind} ile enjekte edilmiş kod tespiti yapılır.

\subsection{Behavioral Analysis ve Family Classification}
Çalışma zamanı sırasında dosya sistemi, Kayıt Defteri ve süreç oluşturma olaylarının izlenmesi. Mutex adları gibi sabit kodlanmış nesnelerin tespiti, \textbf{Mutex} tabanlı aile sınıflandırması için en güvenilir göstergelerdendir.

\subsection{Anti-analysis Evasion Techniques}
Anti-VM ve anti-debugging mekanizmalarının tespiti. Paketlenmiş kötü amaçlı yazılımlarda \textbf{Orijinal Giriş Noktası (OEP)} tespiti ve paket açma rutininin tersine mühendislikle çözülmesi.

\subsection{Automated Malware Analysis ve YARA Rule Development}
YARA kuralları, dosya tabanlı imzalamadan mantıksal ve davranışsal imzalamaya geçişte merkezi bir rol oynar. Bellek dökümü ve dinamik analizden elde edilen \textbf{Mutex} adları ve \textbf{Bellek Paternleri} ile yüksek güvenilirlikli, davranışsal YARA kuralları geliştirilir.

\section{Threat Intelligence Integration ve Operationalization}

\subsection{SIEM ve SOAR Platform Integration}
Tehdit istihbaratının, log korelasyonu, alarm zenginleştirme (\textit{enrichment}) ve otomatik yanıt için SIEM ve SOAR platformlarına API/STIX/TAXII aracılığıyla entegrasyonu.

\subsection{Automated Threat Response ve Playbook Development}
Yüksek güven skoruna sahip IoC'ler temelinde, ağ ve uç nokta seviyesinde hızlı ve kesin engelleme aksiyonlarının \textbf{SOAR Oyun Kitapları} (\textit{Playbooks}) ile otomatikleştirilmesi.

\subsection{Intelligence Sharing Communities ve Trust Groups}
Yapılandırılmış ve güvenilir CTI paylaşımı. \textbf{MISP} (Malware Information Sharing Platform) ile olay merkezli paylaşım ve \textbf{PyMISP} ile otomatik veri çekme ve zenginleştirme.

\subsection{Custom Threat Intelligence Development}
Dahili analizler (sandbox, tehdit avcılığı) sonucunda elde edilen organizasyona özgü göstergelerin, \textbf{STIX 2.1} formatına dönüştürülmesi ve gösterge yaşam döngüsü yönetiminin (\textbf{aging}, \textbf{güven skoru}) otomasyonu.

\subsection{Threat Intelligence ROI Measurement ve Effectiveness}

\subsubsection{Yanlış Pozitif Oranı (FPR) Değerlendirmesi ve İndirgenmesi}
TI kalitesinin, ürettiği \textbf{Yanlış Pozitif (FP)} sayısıyla ters orantılı olduğu anlaşılır. FPR'nin hesaplanması ve IoC validasyonu ile düşürülmesi.

\subsubsection{Operasyonel Etkinlik Metrikleri}
TI'ın değerinin, Ortalama Tespit Süresi (\textbf{MTTD}) ve Ortalama Müdahale Süresi (\textbf{MTTR}) gibi anahtar performans göstergeleri (\textbf{KPI}) üzerindeki ölçülebilir iyileşmelerle kanıtlanması.